\documentclass[12pt,a4paper]{article}
\usepackage[utf8]{inputenc}
\usepackage{geometry}
\usepackage{setspace}
\usepackage{hyperref}
\geometry{margin=1in}
\setstretch{1.2}

\begin{document}

% ---------------------------
% Cover Page
% ---------------------------
\begin{titlepage}
    \centering
    \vspace*{2cm}
    {\LARGE \textbf{Real Estate Price Prediction}} \\[0.5cm]
    {\Large Advanced Apex Project -- Phase 1 Proposal} \\[1.5cm]

    {\large \textbf{Team: The Outliers}} \\[0.5cm]

    \vspace{1cm}
    \begin{tabular}{|p{4cm}|p{10cm}|}
    \hline
    \textbf{Name} & \textbf{Details} \\
    \hline
    Anik Das (Representative) & 
    BITS ID: 2025EM1100026 \newline
    Email: \href{mailto:2025em1100026@bitspilani-digital.edu.in}{2025em1100026@bitspilani-digital.edu.in} \newline
    Contact: +91 9163056329 \\
    \hline
    Adeetya Wadikar & 
    BITS ID: 2025EM1100384 \newline
    Email: \href{mailto:2025em1100384@bitspilani-digital.edu.in}{2025em1100384@bitspilani-digital.edu.in} \newline
    Contact: +91 9653646741 \\
    \hline
    Pallavi Pandey & 
    BITS ID: 2025EM1100084 \newline
    Email: \href{mailto:2025em1100084@bitspilani-digital.edu.in}{2025em1100084@bitspilani-digital.edu.in} \newline
    Contact: +91 7798399666 \\
    \hline
    Tushar Nishane & 
    BITS ID: 2025EM1100306 \newline
    Email: \href{mailto:2025em1100306@bitspilani-digital.edu.in}{2025em1100306@bitspilani-digital.edu.in} \newline
    Contact: +91 9561434710 \\
    \hline
    \end{tabular}

    \vspace{1.5cm}
    \begin{flushleft}
    \textbf{Supervisor:} \underline{\hspace{8cm}} \\[0.5cm]
    \textbf{Course:} Advanced Apex Project 1 (Trimester 1, 2025--26) \\[0.5cm]
    \textbf{Submission:} Phase 1 Proposal \\
    \end{flushleft}

    \vfill
    {\large Date: \today}
\end{titlepage}

% ---------------------------
% Proposal Content
% ---------------------------

\noindent \textbf{Project Title: Real Estate Price Prediction}

\section*{1. Problem Statement}
Accurate real estate price prediction is a major challenge for buyers, sellers, and investors. Traditional valuation methods rely on limited appraisal factors and often fail to capture the diverse influences such as lot size, number of rooms, construction year, location, and neighborhood characteristics. Inaccurate estimates may lead to poor investment decisions, buyer dissatisfaction, or overpricing in competitive markets. We need a data-driven approach to analyze historical housing data and build predictive models that can estimate house sale prices with greater reliability. This solution will also help identify which features most strongly affect property prices.

\section*{2. Business Goal}
The primary goal is to develop a machine learning regression model that predicts property sale prices with strong accuracy and reliability. This will allow potential buyers and investors to make informed decisions, real estate agencies to improve their valuation practices, and stakeholders to understand the market drivers of property prices. By the end of the project, we aim to deliver a working predictive model along with clear insights and visualizations. 

\section*{3. Data Source}
We will use the \textbf{Ames Housing Dataset}, a widely used benchmark dataset for real estate price prediction.  

\textbf{Source Platform:} Kaggle  

\textbf{Full Citation:}  
Shashank Necrothapa. (n.d.). \textit{Ames Housing Dataset} [Data set]. Kaggle. Retrieved October 1, 2025, from \url{https://www.kaggle.com/datasets/shashanknecrothapa/ames-housing-dataset}  

\textbf{Dataset Details:}
\begin{itemize}
    \item Instances: 2,930 observations
    \item Features: 79 attributes (numerical and categorical)
    \item Target: \texttt{SalePrice} (house sale price)
    \item Features describe lot size, rooms, year built, neighborhoods, etc.
\end{itemize}

\section*{4. Tools \& Technologies}
\begin{itemize}
    \item \textbf{Programming Language:} Python
    \item \textbf{Libraries:} Pandas, NumPy, Scikit-learn, Matplotlib, Seaborn
    \item \textbf{Development Environment:} Jupyter Notebook
    \item \textbf{BI Tools:} Tableau or Power BI for final dashboard
    \item \textbf{Version Control:} GitHub
\end{itemize}

\section*{5. Project Workflow}
The project will follow the standard data science lifecycle:

\begin{itemize}
    \item \textbf{Data Acquisition:} Obtain the dataset from Kaggle (direct download method).
    \item \textbf{Preprocessing:} Handle missing values, clean inconsistencies, encode categorical variables.
    \item \textbf{Exploratory Data Analysis (EDA):} Use statistical summaries and visualizations to study patterns.
    \item \textbf{Feature Engineering:} Create new features (e.g., house age, price per sqft) and perform scaling/encoding.
    \item \textbf{Model Building:} Train regression models (Linear Regression, Ridge/Lasso, Random Forest, Gradient Boosting).
    \item \textbf{Model Evaluation:} Evaluate using RMSE, MAE, and R\textsuperscript{2}.
    \item \textbf{Visualization \& Reporting:} Create dashboards and storytelling slides highlighting findings.
\end{itemize}

\section*{6. Data Extraction}
The dataset (Ames Housing) has been acquired using the \textbf{direct download method} from Kaggle.  
We manually downloaded the ZIP file from the Kaggle dataset page and extracted it into the project folder for further preprocessing.  

\textbf{Files:} \texttt{AmesHousing.csv}  
\textbf{Location:} \texttt{./data/}  

This ensures that our dataset remains consistent with the official Kaggle version. A short inspection notebook will be prepared to confirm schema details and enable reproducibility in later phases.

\section*{7. Schema / Data Dictionary}
A schema or data dictionary provides a structured overview of the dataset.  
For Phase~1, we prepared a concise dictionary with 20 key features that are most relevant for modeling house prices.  
The complete dictionary covering all 82 features of the Ames Housing dataset will be maintained separately in an Excel file (\texttt{data\_dictionary.xlsx}) and updated as the project progresses.  

\begin{tabular}{|l|l|p{8cm}|c|}
\hline
\textbf{Feature} & \textbf{Type} & \textbf{Description} & \textbf{PK} \\
\hline
Id & Integer & Unique property identifier & Yes \\
\hline
MSSubClass & Categorical & Type of dwelling (e.g., 20 = 1-Story 1946+ All Styles, 30 = 1-Story 1945) & No \\
\hline
MSZoning & Categorical & General zoning classification (Residential, Commercial, etc.) & No \\
\hline
LotFrontage & Integer & Linear feet of street connected to property & No \\
\hline
LotArea & Integer & Lot size in square feet & No \\
\hline
Street & Categorical & Type of road access (Grvl, Pave) & No \\
\hline
LotShape & Categorical & General shape of property (Reg, IR1, IR2, IR3) & No \\
\hline
Neighborhood & Categorical & Physical location / neighborhood within Ames city & No \\
\hline
Condition1 & Categorical & Proximity to main road or other conditions & No \\
\hline
BldgType & Categorical & Type of dwelling (1Fam, TwnhsE, etc.) & No \\
\hline
HouseStyle & Categorical & House style (1Story, 2Story, etc.) & No \\
\hline
OverallQual & Integer & Overall material and finish quality (scale 1–10) & No \\
\hline
OverallCond & Integer & Overall condition rating (scale 1–10) & No \\
\hline
YearBuilt & Integer & Original construction year & No \\
\hline
YearRemodAdd & Integer & Year of remodel/addition (if any) & No \\
\hline
GrLivArea & Integer & Above-grade (ground) living area square feet & No \\
\hline
FullBath & Integer & Full bathrooms above grade & No \\
\hline
BedroomAbvGr & Integer & Bedrooms above ground level & No \\
\hline
GarageCars & Integer & Garage capacity in number of cars & No \\
\hline
GarageArea & Integer & Garage size in square feet & No \\
\hline
SalePrice & Integer & Property sale price (Target variable) & No \\
\hline
\end{tabular}


\bigskip

\noindent \textbf{Prepared by:} Team The Outliers \\
\textbf{Phase:} Phase 1 — Proposal \\
\textbf{Date:} \today

\end{document}
